\documentclass{IEEEtran}

% ======================
% Packages
% ======================
\usepackage{graphicx}
\usepackage{amsmath}
\usepackage{float}
\usepackage{subcaption}
\title{Fetal Head Circumference Estimation}

\author{
    Pham Cong Duyet--23BI14136\\
    ML in Medicine -- Practical Work 2\\
    University of Science and Technology of Hanoi
}

\begin{document}
\maketitle
\section{Introduction}

Ultrasound is a common medical imaging method because it is safe and easy to use. In pregnancy care, fetal head circumference (HC) is an important measurement to assess fetal growth. Measuring HC manually from ultrasound images can be time-consuming and may lead to different results between observers. In this practical work, I explore the HC18 dataset and build a machine learning model to automatically estimate fetal head circumference using a regression approach.
\section{Dataset Description}
We utilize the HC18 dataset, which consists of \textbf{999 ultrasound scans} for regression analysis. The target head circumference (HC) values range from 50 mm to 350 mm. Despite standardized image contrast, the dataset presents three primary characteristics and challenges:

\begin{itemize}
    \item \textbf{Data Distribution:} The HC values are not uniformly distributed. There are significant peaks (multimodal distribution), with the majority of samples concentrated around 170 mm.
    \item \textbf{Image Quality:} The scans contain varying levels of speckle noise and acoustic shadows, common in ultrasound imaging.
    \item \textbf{Boundary Ambiguity:} High variation in fetal head orientation and low contrast at the skull edges make precise feature extraction difficult.
\end{itemize}

% --- BẮT ĐẦU CHÈN HÌNH SAU DANH SÁCH ---
\begin{figure}[htbp] % 'htbp' để LaTeX tự chọn vị trí tốt nhất (ưu tiên tại chỗ)
    \centering
    
    % Hình A: Ảnh mẫu
    \begin{subfigure}[b]{0.75\linewidth} % Dùng \linewidth để ảnh tự co giãn theo bề ngang cột
        \centering
        \includegraphics[width=\textwidth]{figure/sample.png} 
        \caption{Ultrasound samples with HC annotations}
        \label{fig:samples}
    \end{subfigure}
    
    \vspace{0.4cm} % Khoảng cách giữa 2 hình cho thoáng mắt

    \caption{Visual overview of the HC18 dataset. Illustrates the speckle noise and visual diversity in ultrasound scans.}
    \label{fig:dataset_overview}
\end{figure}
\section{Methodology}
Random Forest Regressor is an ensemble machine learning method that performs regression by averaging the predictions of multiple decision trees. This approach improves robustness and allows the model to capture non-linear relationships between input features and the target variable.

In this work, fetal head circumference estimation was formulated as a regression problem. Each ultrasound image was resized to a fixed resolution and converted into a grayscale representation. The pixel intensities were then flattened into one-dimensional feature vectors and normalized to the range $[0,1]$. These vectors served as input features for the regression model.

The dataset was split into training and validation subsets using an 80/20 ratio with a fixed random seed to ensure reproducibility. A Random Forest Regressor was implemented using the \textit{scikit-learn} library. The baseline configuration used 100 decision trees without depth limitation, which provided a good balance between accuracy and computational cost. Model performance was evaluated using the Mean Absolute Error (MAE)

To assess the influence of hyperparameters on performance, additional experiments were conducted by varying the number of trees (\texttt{n\_estimators} = 50, 100, 200) and the maximum tree depth (\texttt{max\_depth} = Unlimited, 10, 20). These experiments allowed us to analyze model stability and sensitivity to hyperparameter changes.
\section{Results}

\subsection{ Results}

The performance of the proposed Random Forest Regressor was evaluated on the validation set using the Mean Absolute Error (MAE) metric. The baseline configuration of the model, using 100 decision trees without depth limitation, achieved an MAE of \textbf{27.15 mm}. This result serves as the reference performance for further comparison.

To analyze the influence of model hyperparameters, additional experiments were conducted by varying the number of trees (\texttt{n\_estimators}) and the maximum depth of each tree (\texttt{max\_depth}). The quantitative results of these experiments are summarized in Table~\ref{tab:rf_results}.

The results indicate that increasing the number of trees generally leads to a slight improvement in performance. The best result was obtained with \texttt{n\_estimators = 200} and \texttt{max\_depth = Unlimited}, achieving an MAE of \textbf{27.13 mm}. However, the improvement compared to the baseline configuration is marginal, suggesting stable model behavior across different hyperparameter settings. Notably, removing depth constraints allows the trees to grow fully, which appears beneficial for capturing the underlying non-linear relationships in the data.
\begin{table}[h!]
\centering
\caption{Random Forest regression results with different hyperparameters}
\label{tab:rf_results}
\begin{tabular}{ccc}
\hline
n\_estimators & max\_depth & MAE (mm) \\
\hline
50  & Unlimited & 27.49 \\
50  & 10   & 27.70 \\
50  & 20   & 27.65 \\
100 & Unlimited & 27.15 \\
100 & 10   & 27.33 \\
100 & 20   & 27.18 \\
200 & Unlimited & \textbf{27.13} \\
200 & 10   & 27.38 \\
200 & 20   & 27.25 \\
\hline
\end{tabular}
\end{table}
\subsection{Comparison with Leaderboard}

The obtained results were compared with the HC18 Grand Challenge leaderboard. Top-ranking methods report mean absolute errors below 1 mm, which are achieved using advanced deep learning architectures such as U-Net variants and nnU-Net, combined with pixel-wise anatomical segmentation and extensive training strategies.

In contrast, the proposed approach relies on a traditional machine learning model operating on flattened image features, without explicit anatomical segmentation or deep feature extraction.This comparison highlights the trade-off between model complexity and performance in medical image regression tasks.

\section{Conclusion}
Experimental results show that the Random Forest model provides stable performance across different hyperparameter configurations, achieving a minimum MAE of 27.13 mm on the validation set. Although this performance is lower than state-of-the-art methods reported on the HC18 Grand Challenge leaderboard, the comparison highlights the impact of model complexity on prediction accuracy.

Overall, this work demonstrates that traditional machine learning techniques can serve as effective baseline models for medical image regression tasks. Future improvements may include incorporating anatomical feature extraction, segmentation-based approaches, or deep learning models to further enhance prediction accuracy.

\end{document}